\chapter{True Is Not False, at the Boundary}
\label{chap:true-not-false-at-the-boundary}

Early in the book the construction did something that should have been alarming and was instead
deliberate: it collapsed the distinction between true and false. The bare interior apparatus, identifying
any two readings it could not observationally tell apart, could not tell a true reading from a false one
either --- the quotient that gave the construction its spare foundations also flattened truth and
falsehood into one. The book has run ever since on an apparatus that does not, inside, distinguish the
two. This chapter restores the distinction, and it restores it exactly where the outside readings have
been living: at the boundary. There the signal is on or off, the two are provably not equal, and for a
distinguished sentence the true reading and the false reading at last differ. Having recovered the
distinction, the chapter asks how much of it the reader wants, offering a ladder of resolution from a
single bit up to full inference and letting the reader stop wherever they choose --- because, as the last
section shows, one distinguishable bit is already enough.

\section{The distinction restored}
\label{se:the-distinction-restored}

The collapse was real and the restoration must be equally real to mean anything. Inside, the
construction's quotient identifies observationally indistinguishable readings, and under it a true fact
and a false fact, if nothing the apparatus can observe separates them, are the same fact. This is the
flattening the opening chapters performed and the rest of the book accepted. At the boundary the
flattening lifts. The boundary signal takes one of two values, on or off, and the construction proves
they are not equal --- on is not off, off is not on, the inequality holding as a checked fact rather than
a convention. It gives falsehood its own standing as well: a genuine false fact, distinct from the true,
not a true fact in disguise. And it proves the consequence that matters: true is not false. For a
distinguished sentence --- one of the self-referential sentences of the last chapter --- the reading that
takes it true and the reading that takes it false are different readings, separated at the boundary by
the on/off signal that the interior could not supply.

It is worth being careful about what was and was not collapsed, because the restoration depends on the
distinction being recoverable rather than merely re-declared. The interior did not abolish truth and
falsehood; it declined to distinguish them where it could observe no difference. A true reading and a
false reading the apparatus cannot tell apart are, to the apparatus, the same reading --- not because
truth has been denied but because the apparatus has no instrument fine enough to separate them. What the
boundary supplies is exactly that missing instrument. The on/off signal is something the interior could
not produce and the boundary can, and with it the two readings that were indistinguishable inside become
distinguishable outside. So the restoration is not a reversal of the collapse but its completion: the
distinction was always there in principle, waiting for an instrument fine enough to read it, and the
boundary is where that instrument finally exists.

This is the law of non-contradiction, recovered. The experiment under the names of Aristotle and De
Morgan marks what is at stake: that a thing and its negation cannot both hold, the oldest constraint in
logic, which the interior's quotient had suspended and the boundary now reinstates. The experiment on the
positron threshold supplies the physical echo --- the restored distinction is the same one that separates
the two orientations, the plus and the minus that the trace of three chapters ago could read precisely
because, at the boundary, they are not the same. The construction spent the distinction at the start to
buy its spare foundations; it earns it back here, at the edge, as a theorem.

The recovery of non-contradiction at the boundary has a quiet structural elegance worth pausing on. The
interior, having flattened true and false, could not state the law that they exclude each other --- there
was nothing to exclude, the two being one. The boundary, by separating them, makes the law statable
again, and the construction can finally write down that the on signal and the off signal cannot both
hold, the exclusion the interior could not express. What is striking is that the law returns not as an
axiom imposed but as a consequence of the boundary's having two distinct values: non-contradiction is not
added to the construction, it falls out of on being unequal to off. The oldest law of logic, suspended at
the start for the sake of spare foundations, is recovered at the end as a theorem about a two-valued
signal --- proof that the construction never needed to assume it, only to build far enough to derive it.

\section{The ladder, and how much you want}
\label{se:the-ladder-and-how-much-you-want}

A restored distinction raises a question the construction is unusually willing to hand to the reader: how
much resolution do you want? Between the bare boundary bit and a fully reasoned verdict there is a
ladder, and the construction builds it explicitly. Its bottom rung is the binary signal, the single
on/off bit. Its top rung is the level the construction calls inferred --- the closure calculus, where
facts are not merely read but reasoned to their consequences --- and the construction proves the ladder
closes there, inferred sealed as the last rung, with definite stages in between carrying the bit upward
through richer and richer readings. The whole ladder, binary at the bottom and inferred at the top, is a
finite tower the construction can exhibit.

The willingness to hand the reader the choice of depth is a stance, not a shrug, and it is worth seeing
why the construction can afford it. A theory that hid its lower rungs --- that insisted you take its fully
reasoned verdict or nothing --- would be asking for trust the construction has spent the whole book
refusing to ask for. By exposing the ladder, the construction lets a reader who wants only the bare
distinction take the bottom bit and verify it directly, and lets a reader who wants the full reasoning
climb to inference, and treats both as legitimate readings of the same fact at different resolutions. This
is the reciprocal accounting of the orientation chapter made operational: a reading at any rung is a
reading of the whole, taken at that depth, with the unclimbed rungs accounted for as available rather than
hidden. The reader who stops low has not been deceived about what lies higher; the higher rungs are on the
books, offered, declined for now.

What is unusual is that the reader chooses where to stop. The construction provides a recursion-stop
rule --- a principle that halts the climb at whatever rung the reader asks for --- so that one may take
the bare bit and go no further, or climb to a richer reading, or run all the way to inference, and each
choice is legitimate. A formal theorem authorizes going deeper whenever more is wanted, so the depth is
not capped by the construction but offered as the reader's to set. The two acts run the ladder: the
reader funges the stages into the resolution they want, pooling the rungs up to their chosen depth, and
tanges the stop, selecting by the recursion-stop rule the exact rung to halt on. Governing the bit at
the bottom is the duality the experiment on De Morgan names: the on and the off, the meet and the join,
related by the same negation-swapping symmetry that turns an and into an or, so that the single bit
already carries the logical structure the higher rungs unfold. The construction does not dictate a
resolution; it builds the ladder and lets the reader say how much they want.

The De Morgan duality at the bottom rung deserves one more word, because it is what makes a single bit
logically complete rather than merely binary. De Morgan's laws say that negation turns conjunction into
disjunction and back --- that to deny both of two things hold is to affirm that one of them fails. On the
boundary's two-valued signal this duality is exact: the on and the off, the meet and the join, are
swapped by negation, so that whichever of the two the reader takes as primary, the other and their
combinations are fixed. This is why the bottom rung is not a mere flag but a foothold for logic. A reader
who funges the two values into their conjunction and one who funges them into their disjunction are
related by a negation the construction can run, and either can tange out the value they need. The single
bit, governed by De Morgan, already holds the and, the or, and the not --- the whole of the propositional
logic the higher rungs will use, packed into one distinguishable signal. That a reader stopping at the
bottom rung loses none of that logic is the deepest sense in which one bit is enough: the bit is not below
logic but already logical, the smallest object on which conjunction, disjunction, and negation all act.

\section{One bit is enough}
\label{se:one-bit-is-enough}

The chapter's last move is to show that the bottom rung already suffices. The construction proves that
one distinguishable bit is enough: the single on/off signal at the boundary carries the whole restored
distinction, and a reader who climbs no higher than that bit has, nonetheless, everything needed to tell
true from false. The distinguished sentence is distinguishable --- its true and false readings differ ---
and they differ already at the one bit, so the higher rungs add resolution but not necessity. Whatever
more the ladder offers is refinement; the distinction itself lives in the bit.

The cleanest instance is the orthogonality the book has carried since the split: the gauge sector and
the geometric sector are perpendicular, and the construction now states it in the chapter's own terms ---
the QED side and the GR side are orthogonal, distinguishable by a single bit. Two whole sectors of the
construction's physics, the charge-bearing and the geometric, are told apart by one binary distinction,
on or off, true or false. This is the economy the whole part has been building toward. The outside reading
compressed an interior physics to a bit; the obstruction was a bit; the undecidable sentence turned on a
bit; and now the restored distinction, the recovered law of non-contradiction, the separation of two
physical sectors --- all of it comes down to whether a single signal is on or off. One bit, distinguished
and distinguishable, is enough.

That so much reduces to one bit is not a trivialization but a result, and the difference is worth marking.
It would be a trivialization to say the construction's physics is \emph{merely} a bit, that nothing of its
richness survives the compression. That is not the claim. The interior richness is real and stays real ---
the charge, the sign, the invariant, the whole structure the inside reads. What the chapter proves is
narrower and stranger: that the single question an outside observer needs answered --- is the distinction
there or not --- is settled by one bit, however rich the interior that bit summarizes. The bit is not the
physics; it is the physics' answer to a yes-or-no question, and the result is that the yes-or-no question
has a one-bit answer. A reader at the boundary gets a bit not because there is only a bit to get but
because a bit is the exact size of the question they are entitled to ask from outside.

\section{What is demonstrated, and what is assumed}
\label{se:demonstrated-assumed-ch28}

The accounting is almost all demonstration, and what little it assumes is the grant the book has carried
throughout.

What is demonstrated is the restoration, the ladder, and the sufficiency. On is proved not equal to off,
falsehood is given its own standing, and true is proved not equal to false at the boundary, with the
distinguished sentence's two readings separated. The decidability ladder from the binary bit to inferred
is built and proved to close at the top, with a recursion-stop rule that lets the reader halt at any
rung and a theorem authorizing more depth on demand. And one distinguishable bit is proved enough to
carry the distinction, with the gauge and geometric sectors shown orthogonal --- told apart by that one
bit. Each is a checked statement about finite readings.

\begin{quote}
\emph{A note on what was spent and earned.} There is no new modeling bridge here, and no new grant; what
deserves marking is the symmetry between this chapter and the opening ones. At the start the construction
spent the true/false distinction deliberately, quotienting it away to buy foundations of remarkable
spareness --- the two axioms of the sixth-problem chapter. That spending was honest and it was real: the
interior genuinely cannot tell true from false. This chapter shows the distinction was not destroyed but
deferred, recoverable at the boundary as a theorem rather than an assumption. The construction did not
smuggle truth back in; it earned it back, at the edge, from the same on/off obstruction the outside
readings have been built on. What was spent inside is repaid outside, and the books balance: a
construction that flattened truth to buy its foundations, and bought back, at its boundary, exactly the
one bit of distinction that flattening cost.
\end{quote}

What is assumed, then, is the standing grant alone, with the restoration proved rather than posited. The
outside readings are now complete: a trace, an obstruction, an undecidable sentence, and a restored
distinction, all resolved to bits at the boundary, all answerable to the reader's chosen depth. The final
part turns from what an outside observer reads to the relations among observers themselves --- relativity,
the second great orientation reversal, and the reader's own standing in the construction they have been
reading.
